\documentclass[12pt]{article}
\usepackage{amsmath, amssymb}
\usepackage[margin=1in]{geometry}
\setlength{\parskip}{6pt}
\title{QUBO Formulation: Weighted Set Cover Problem}
\date{}
\begin{document}
\maketitle

\subsection*{Weighted Set Cover Problem}

\paragraph{Problem Statement.}
Given a finite universe $U$ of $N$ elements and a collection $\mathcal{S} = \{S_0, S_1, \dots, S_{M-1}\}$ of $M$ subsets of $U$, where each subset $S_j$ is associated with a non-negative cost $c_j$, the goal is to select a subcollection of subsets such that every element in $U$ is contained in at least one selected subset, while minimizing the total cost of the selected subsets.

\paragraph{Decision Variables.}
For each subset $j \in \{0, 1, \dots, M-1\}$, introduce a binary variable $x_j \in \{0, 1\}$, where $x_j = 1$ indicates that subset $S_j$ is selected, and $x_j = 0$ otherwise. The decision vector is $\mathbf{x} = (x_0, x_1, \dots, x_{M-1})^\top \in \{0, 1\}^M$.

\paragraph{QUBO Derivation.}
\textbf{Step 1 --- Original Objective.} The problem requires minimizing the total cost of selected subsets:
\begin{align}
\min_{\mathbf{x} \in \{0,1\}^M} \quad C(\mathbf{x}) = \sum_{j=0}^{M-1} c_j x_j.
\end{align}
In the QUBO framework, we minimize a single quadratic Hamiltonian $H(\mathbf{x})$ that combines the objective with constraint penalties.

\textbf{Step 2 --- Constraint Formulation.} For each element $i \in \{0, 1, \dots, N-1\}$, let $I_i = \{j \in \{0, \dots, M-1\} : i \in S_j\}$ denote the set of indices of subsets that cover element $i$. The coverage requirement is that each element must be covered at least once:
\begin{align}
\sum_{j \in I_i} x_j \ge 1, \quad \forall i \in \{0, 1, \dots, N-1\}.
\end{align}

\textbf{Step 3 --- Penalty Term Expansion.} Inequality constraints are enforced by adding quadratic penalty terms to the objective. For each element $i$, we introduce a penalty weight $\lambda_i > 0$ and define the penalty function:
\begin{align}
P_i(\mathbf{x}) = \lambda_i \left(1 - \sum_{j \in I_i} x_j\right)^2.
\end{align}
Expanding the square and applying the idempotent property $x_j^2 = x_j$ for binary variables yields:
\begin{align}
\left(1 - \sum_{j \in I_i} x_j\right)^2 &= 1 - 2\sum_{j \in I_i} x_j + \left(\sum_{j \in I_i} x_j\right)^2 \\
&= 1 - 2\sum_{j \in I_i} x_j + \sum_{j \in I_i} x_j^2 + \sum_{\substack{j,k \in I_i \\ j \neq k}} x_j x_k \\
&= 1 - \sum_{j \in I_i} x_j + 2\sum_{\substack{j,k \in I_i \\ j < k}} x_j x_k.
\end{align}
Thus, the expanded penalty contribution for element $i$ is:
\begin{align}
P_i(\mathbf{x}) = \lambda_i - \lambda_i \sum_{j \in I_i} x_j + 2\lambda_i \sum_{\substack{j,k \in I_i \\ j < k}} x_j x_k.
\end{align}

\textbf{Step 4 --- Combined Hamiltonian.} Summing the objective and all penalty terms gives the full QUBO Hamiltonian:
\begin{align}
H(\mathbf{x}) = \sum_{j=0}^{M-1} c_j x_j + \sum_{i=0}^{N-1} \left( \lambda_i - \lambda_i \sum_{j \in I_i} x_j + 2\lambda_i \sum_{\substack{j,k \in I_i \\ j < k}} x_j x_k \right).
\end{align}
Collecting coefficients for linear and quadratic terms:
\begin{align}
H(\mathbf{x}) = \sum_{j=0}^{M-1} \left( c_j - \sum_{i: j \in I_i} \lambda_i \right) x_j + \sum_{0 \le j < k \le M-1} \left( 2 \sum_{i: j,k \in I_i} \lambda_i \right) x_j x_k + \sum_{i=0}^{N-1} \lambda_i.
\end{align}

\textbf{Step 5 --- Q Matrix and Offset.} Reading off the coefficients from the expanded Hamiltonian, the upper-triangular QUBO matrix entries $Q_{j,k}$ and constant offset $c$ are given in closed form as:
\begin{align}
Q_{j,j} &= c_j - \sum_{i: j \in I_i} \lambda_i, \quad \forall j \in \{0, \dots, M-1\}, \\
Q_{j,k} &= 2 \sum_{i: j,k \in I_i} \lambda_i, \quad \forall 0 \le j < k \le M-1, \\
c &= \sum_{i=0}^{N-1} \lambda_i.
\end{align}

\paragraph{Penalty Weight Justification.}
The penalty weights $\lambda_i$ must be chosen large enough to ensure that violating any coverage constraint increases the total energy more than any possible reduction in the objective function. The maximum possible cost reduction achievable by omitting subsets is bounded by the sum of all subset costs, $\sum_{j=0}^{M-1} c_j$. A sufficient condition to guarantee that the global minimum lies in the feasible region is to set $\lambda_i = \lambda > \sum_{j=0}^{M-1} c_j$ for all $i$. Under this condition, the penalty incurred by leaving any element uncovered strictly exceeds the maximum objective gain from any subset selection, thereby enforcing feasibility.

\paragraph{Correctness Argument.}
Minimizing $H(\mathbf{x})$ over $\mathbf{x} \in \{0,1\}^M$ recovers the optimal set cover because the penalty terms are non-negative and vanish exactly when all coverage constraints are satisfied. For any feasible solution, $H(\mathbf{x})$ reduces to the original cost function $C(\mathbf{x})$. For any infeasible solution, at least one penalty term is strictly positive. By choosing $\lambda$ sufficiently large, the energy of any infeasible configuration exceeds the energy of the optimal feasible configuration. Consequently, the global minimum of $H(\mathbf{x})$ must occur at a feasible point, and among feasible points, $H(\mathbf{x})$ is identical to the original objective, ensuring the global minimum corresponds to the minimum-cost set cover.

\paragraph{Illustrative Example.}
Consider a concrete instance with $M=3$ subsets and $N=3$ elements. Let the subset costs be $c_0 = 10$, $c_1 = 20$, and $c_2 = 15$. The coverage structure is defined by $I_0 = \{0, 2\}$, $I_1 = \{0, 1\}$, and $I_2 = \{1, 2\}$. We set a uniform penalty weight $\lambda = 100$. Applying the general formulas from Step 5:
\begin{align}
Q_{0,0} &= 10 - (100 + 100) = -190, \\
Q_{1,1} &= 20 - (100 + 100) = -180, \\
Q_{2,2} &= 15 - (100 + 100) = -185, \\
Q_{0,1} &= 2(100) = 200, \\
Q_{0,2} &= 2(100) = 200, \\
Q_{1,2} &= 2(100) = 200, \\
c &= 100 + 100 + 100 = 300.
\end{align}
The resulting Hamiltonian is:
\begin{align}
H(x_0, x_1, x_2) &= -190x_0 - 180x_1 - 185x_2 + 200x_0x_1 + 200x_0x_2 + 200x_1x_2 + 300.
\end{align}
Evaluating $H$ over all $2^3 = 8$ binary configurations, the feasible solutions are $(1,0,1)$, $(0,1,1)$, $(1,1,0)$, and $(1,1,1)$. Their energies are $H(1,0,1) = 215$, $H(0,1,1) = 215$, $H(1,1,0) = 220$, and $H(1,1,1) = 255$. The minimum energy is $215$, achieved by $\mathbf{x}^* = (1,0,1)$ and $\mathbf{x}^* = (0,1,1)$, corresponding to selecting subsets $\{S_0, S_2\}$ or $\{S_1, S_2\}$ with total cost $25$. This matches the optimal solution of the original weighted set cover instance.

\end{document}